%%%%%%%%%%%%%%%%%%%%%%%%%%%%%%%%%%%%%%%%%
% FRI NLP Report
%%%%%%%%%%%%%%%%%%%%%%%%%%%%%%%%%%%%%%%%%

\documentclass[fleqn,moreauthors,10pt]{ds_report}
\usepackage[english]{babel}
\usepackage{hyperref}
\graphicspath{{fig/}}

\JournalInfo{FRI Natural Language Processing Course 2026}
\Archive{Project report}

\PaperTitle{Domain-Specific Fitness NLP Assistant Using Retrieval-Augmented Generation}

\Authors{Domen Pahole, Luka Rizman, Nina Trivić}

\affiliation{\textit{Advisors: Aleš Žagar}}

\Keywords{Natural Language Processing, Fitness NLP, Retrieval-Augmented Generation, Semantic Retrieval, Domain Adaptation}
\newcommand{\keywordname}{Keywords}

\Abstract{
This project explores the development of a domain-specific natural language processing assistant focused on fitness-related textual knowledge. The proposed system leverages structured exercise descriptions and training program data to enable semantic retrieval, recommendation, and contextual response generation.
}

\begin{document}
\flushbottom
\maketitle
\thispagestyle{empty}

\section*{Introduction}

The increasing availability of digital fitness content has resulted in large textual corpora describing exercise routines, nutrition plans, and training recommendations. While general-purpose language models demonstrate strong linguistic capabilities, they often lack domain grounding required for reliable reasoning in specialized areas such as fitness and health.

This project aims to develop a domain-specific NLP assistant capable of understanding user queries related to fitness training and providing grounded responses based on curated textual knowledge. The assistant will rely on structured corpora containing exercise descriptions and training plans. Retrieval-augmented generation techniques will be explored to improve factual reliability and contextual relevance.

\section*{Datasets}

The project utilizes multiple publicly available datasets to construct a fitness-related textual corpus.

\begin{itemize}
\item \href{https://www.kaggle.com/datasets/exercisedb/fitness-exercises-dataset}{ExerciseDB Dataset}
\item \href{https://www.kaggle.com/datasets/ambarishdeb/gym-exercises-dataset}{Gym Exercise Dataset}
\end{itemize}

\subsection*{Nutrition Knowledge Corpora}

\begin{itemize}
\item \href{https://www.kaggle.com/datasets/utsavdey1410/food-nutrition-dataset}{Food Nutrition Dataset}
\item \href{https://www.kaggle.com/datasets/adilshamim8/daily-food-and-nutrition-dataset}{Daily Food and Nutrition Dataset}
\item \href{https://fdc.nal.usda.gov/download-datasets}{USDA FoodData Central Dataset}
\end{itemize}

\subsection*{Fitness User Context Datasets}

\begin{itemize}
\item \href{https://www.kaggle.com/datasets/valakhorasani/gym-members-exercise-dataset}{Gym Members Exercise Dataset}
\item \href{https://www.kaggle.com/datasets/nadeemajeedch/fitness-tracker-dataset}{Fitness Tracker Dataset}
\end{itemize}

\section*{Initial Project Idea}

The project proposes the development of a fitness knowledge assistant capable of responding to natural language queries. The assistant will retrieve relevant exercise descriptions and generate context-aware responses.

Potential applications include:

\begin{itemize}
\item Exercise recommendation based on training goals
\item Semantic similarity search between exercises
\item Summarization of workout instructions
\item Interpretation of training queries
\end{itemize}

\section*{Results}

At this stage of the project, we have collected the datasets and reviewed them, and we have defined the first preprocessing steps. The data comes from different sources and includes exercise descriptions, nutrition information, and workout plans. Since the datasets are in different formats, we need to clean, standardize, and extract the text before using them further.

The exercise data includes structured information such as exercise name, target muscle group, equipment, and instructions. The nutrition data includes food names, nutritional values, and serving sizes. This information can be used to build a text-based knowledge base for search and recommendations.

As a first step, we will build a simple baseline pipeline to test how well retrieval works. This will include creating text embeddings using transformer-based models and storing them in a vector database. We will then test if the system can correctly retrieve relevant exercises based on simple queries, such as searching by muscle group or workout goals.

\section*{Related Work}

Recent research shows that artificial intelligence methods can support fitness recommendation and exercise guidance. Yang et al. \cite{Yang2023} proposed a physical exercise recommendation system that combines statistical methods and natural language processing, showing that NLP can be useful for personalized fitness assistance.

For domain-specific question answering, Retrieval-Augmented Generation (RAG) is a useful approach. Lewis et al. \cite{Lewis2020} introduced RAG as a method that combines a language model with document retrieval, allowing answers to be generated from external knowledge. This is relevant for our project, since a fitness assistant should respond using exercise and nutrition data instead of unsupported general knowledge.

\bibliographystyle{unsrt}
\bibliography{report}

\end{document}